% ------------------------------------------------------------------
% SSUS working-paper format — end-of-preamble definitions
% hikmah carried much of this in its before-body partial; on quarto
% >= 1.10 preamble content splices via before-title.tex.
% ------------------------------------------------------------------

% -----------------
% Brand fonts
% -----------------
% IBM Plex Serif 400/600 for what is read at length (body, abstract,
% footnotes, references, tables) and IBM Plex Math, its companion, for
% mathematics; Libre Franklin 400/800 for the masthead, headings, and
% captions; IBM Plex Mono for code, data, and the UC IRVINE line (brand
% README, *Type*). In the sans, bold maps to ExtraBold (800), the brand's
% display weight; in the serif, to SemiBold. Plex Serif and Plex Math
% are TeX Live's v3.001 / v1.000 files: Google Fonts' Plex Serif (v2.6)
% carries no Greek. Plex Serif's figures are tabular and lining by
% default, so table columns align; Libre Franklin has no tabular figures.
% The fonts are read in place from the extension's assets/fonts/, whose
% path masthead.lua supplies, so the extension works at any install depth
% and nothing is copied beside the input. Declared here rather than in
% include-in-header.tex, which is not a template; nothing in that header
% sets a font-relative length at preamble time.
\usepackage{fontspec}
\setmainfont{IBMPlexSerif}[
  Path        = $ssus-masthead-fonts$,
  Extension   = .otf,
  UprightFont = *-Regular,
  ItalicFont  = *-Italic,
  BoldFont    = *-SemiBold,
  BoldItalicFont = *-SemiBoldItalic,
]
% pandoc's template loads unicode-math before this header
\setmathfont{IBMPlexMath-Regular.otf}[Path = $ssus-masthead-fonts$]
\setsansfont{LibreFranklin}[
  Path        = $ssus-masthead-fonts$,
  Extension   = .ttf,
  UprightFont = *-Regular,
  ItalicFont  = *-Italic,
  BoldFont    = *-ExtraBold,
  BoldItalicFont = *-ExtraBoldItalic,
]
% IBM Plex Mono vendored beside Libre Franklin (OFL); SemiBold is its 600
\setmonofont{IBMPlexMono}[
  Path        = $ssus-masthead-fonts$,
  Extension   = .ttf,
  UprightFont = *-Regular,
  BoldFont    = *-SemiBold,
]
% Letter-spaced mono, for the masthead series line (mock: .13em)
\newfontfamily\ssusmonoLS{IBMPlexMono}[
  Path        = $ssus-masthead-fonts$,
  Extension   = .ttf,
  UprightFont = *-Regular,
  BoldFont    = *-SemiBold,
  LetterSpace = 15,
]

% Epigraph styling (hikmah core)
\renewcommand{\epigraphsize}{\sffamily\footnotesize}
\setlength{\epigraphrule}{0em}
\setlength{\beforeepigraphskip}{-2em}
\setlength{\afterepigraphskip}{1em}

$if(biblatex)$
% ----------
% BibLaTeX
% ----------
$if(biblatex-chicago)$
\usepackage[$if(biblio-style)$$biblio-style$,$endif$$for(biblatexoptions)$$biblatexoptions$$sep$,$endfor$]{biblatex-chicago}
$else$
\usepackage[$if(biblio-style)$style=$biblio-style$,$endif$$for(biblatexoptions)$$biblatexoptions$$sep$,$endfor$]{biblatex}
$endif$

\setlength\bibitemsep{0pt}  % No space between bib entries
\renewcommand*{\bibfont}{\footnotesize}  % Use smaller font
\setlength\bibhang{\parindent}  % Match document indentation

% Fix biblatex's odd preference for using In: by default.
\renewbibmacro{in:}{%
  \ifentrytype{article}{}{%
  \printtext{\bibstring{}\intitlepunct}}}

\addbibresource{$bibliography$}
$endif$

% Masthead options (from the document's metadata)
\newif\ifssusmark
\ssusmarktrue
$if(masthead-typographic)$
\ssusmarkfalse
$endif$
$if(masthead-mark-file)$
\def\ssusmarkfile{$masthead-mark-file$}
$else$
\def\ssusmarkfile{$ssus-masthead-assets$ssus-mark-light.pdf}  % in place, by masthead.lua
$endif$
$if(masthead-series)$
\def\ssusseries{$masthead-series$}
$else$
\def\ssusseries{SSUS working paper series}
$endif$
$if(date)$
\def\ssusdate{$date$}
$else$
\def\ssusdate{\today}
$endif$